\documentclass{scdpg}
\begin{document}
\scBookLanguage{de}
\begin{scAbstract}
\scLanguage{en}
\scTitle{Deprotonation of quinacridone on Ag(100) - an NIXSW/XPS study of the phase transition}
\scAuthor{*}{Lucas}{Hirschfeld}{1}
\scAuthor{*}{Morris E. L.}{Mühlpointner}{1}
\scAuthor{*}{Anna J.}{Kny}{1}
\scAuthor{*}{Sergey}{Subach}{2}
\scAuthor{*}{Moritz}{Sokolowski}{1}
\scAffiliation{1}{Clausius Institute for Physical and Theoretical Chemistry, University of Bonn, Germany}
\scAffiliation{2}{Peter Grünberg Institut, Forschungszentrum Jülich GmbH, Germany}
\scBeginText

Quinacridone (QA) forms two different phases on the Ag(100) surface. When QA is deposited at room temperature, the formation of the $\alpha$-phase with a point-on-line (\textit{pol}) structure is observed. After annealing, the commensurate $\beta$-phase is irreversibly formed~[1]. The mechanism, involving the loss of intermolecular hydrogen bonds, was not disclosed so far. Thus we have performed a detailed Normal Incidence X-ray Standing Waves (NIXSW)/XPS study.

The N1s XPS spectra exhibit a second peak shifted by 2.2~eV to lower binding energies when going from the $\alpha$- to the $\beta$-phase. This reveals that the irreversible reaction from the $\alpha$- to the $\beta$-phase is explained by a deprotonation of 50~\% of the amine groups. The deprotonated N-atoms form stronger bonds to the Ag(100) surface, indicated by a drastic approach of the respective N adsorption height from 4.3 Å in the $\alpha$-phase to 3.4 Å in the $\beta$-phase. These results show that the reaction is driven by the energy gain related to N/Ag interfacial bonds, which overcompensate the loss of intermolecular hydrogen bonds.

Support by Diamond Light Source is acknowledged.

[1] JPCC 124.45 (2020), 24861–24873.


\scEndText
\scConference{Dresden 2026}
\scPart{O}
\scContributionType{Poster}
\scTopic{Organic molecules on inorganic substrates: electronic, optical and other properties}
\scKeywords{NIXSW; XPS; Quinacridone; Ag(100); Deprotonation}
\scEmail{hirschfeld@pc.uni-bonn.de}
\scCountry{Germany}
\end{scAbstract}
\end{document}
